% TOP_PROBLEM: {"problemId": "problem.sidorenko-conjecture-on-bipartite-homomorphism-densities", "canonicalTitle": "Sidorenko Conjecture on Bipartite Homomorphism Densities", "releaseRank": 265, "primaryDomain": "combinatorics_discrete_geometry", "primaryDomainLabel": "Combinatorics & discrete geometry", "displayStatus": "open", "releaseStatus": "open"}
% !TeX program = lualatex
% Definition and sources only; this file is not an LLM solution attempt.
\documentclass[11pt]{article}
\usepackage[margin=25mm]{geometry}
\usepackage{fontspec}
\setmainfont{Latin Modern Roman}
\usepackage{amsmath,amssymb,mathtools}
\usepackage{unicode-math}
\setmathfont{Latin Modern Math}
\usepackage{xurl,hyperref}
\hypersetup{colorlinks=true,urlcolor=blue}
\setcounter{secnumdepth}{0}
\setlength{\parindent}{0pt}
\setlength{\parskip}{5pt}
\setlength{\emergencystretch}{3em}
\begin{document}
\section*{\texorpdfstring{Sidorenko Conjecture on Bipartite Homomorphism Densities}{Sidorenko Conjecture on Bipartite Homomorphism Densities}}\label{265}
\subsection{Definitions and mathematical statement}
A graphon is a symmetric measurable function $W:[0,1]^2\to[0,1]$. For a finite simple graph $H=(V,E)$, define its homomorphism density by
\[
 t(H,W)=\int_{[0,1]^{V}}\prod_{\{u,v\}\in E}W(x_u,x_v)\prod_{u\in V}{\,\mathrm{d}} x_u.
\]
Sidorenko's assertion is
$t(H,W)\ge t(K_2,W)^{|E(H)|}$ for every bipartite $H$ and every graphon. For the constant graphon $W\equiv p$ there is equality, explaining the quasirandom comparison. Homomorphisms may identify vertices; replacing their density by the density of induced copies or injective embeddings changes the finite-object statement.
\par\smallskip\textit{Formulation and concept references:} \href{https://arxiv.org/pdf/2108.06599}{[S1]}.

\subsection{Short English statement}
Among graph limits with a fixed edge density, does the constant graphon minimize the homomorphism density of every bipartite graph?
\subsection{Sources}
\begin{itemize}
\item[S1]D. Conlon, J. Lee, A. Sidorenko, 'Biregularity in Sidorenko's Conjecture', arXiv:2108.06599 (2021).
\url{https://arxiv.org/pdf/2108.06599}.
\textit{Formulation source.}
\end{itemize}
\end{document}
